\section{第\ref{sec:nora}章　\nt{NORA}}

\nt{NORA}系の構成、その二つのモード、瞳孔との関係（\nt{NORA}経由だけではない）、信号対背景比の上昇、行動における逆U字は直接測定されている。乗算的なゲイン$g$はモデルである。ゲインが選択のモードを切り替え、逆温度として働くことは本章の導出である。ゲイン調整の利点は本書のモデルによる計算である。「\nt{NORA}は探索のつまみ」と「\nt{NORA}は割り込み」は仮説である。

{\footnotesize
\setlength{\tabcolsep}{3pt}\renewcommand{\arraystretch}{1.15}
\begin{longtable}{>{\raggedright}p{0.5cm}>{\raggedright}p{4.0cm}>{\raggedright}p{5.6cm}>{\raggedright}p{2.3cm}>{\raggedright\arraybackslash}p{1.7cm}}
\toprule
No. & 主張 & 実験と測定内容 & 出典 & 状態 \\
\midrule\endfirsthead
\toprule
No. & 主張 & 実験と測定内容 & 出典 & 状態 \\
\midrule\endhead
1 & ヒトの\nt{NORA}ニューロンは数万個で、ほぼすべてが一つの集団にある & ヒト脳幹の連続切片、チロシン水酸化酵素染色。青斑核に$53\,900$個、隣接する亜核にさらに$6\,260$個。 & \cite{Baker1989LC} & 測定済み \\
2 & 出力は脳のほぼ全体に広がるが、集団の部分ごとにある程度異なる領域を担当する & マウスでのウイルス遺伝学的標識。青斑核の\nt{NORA}ニューロンは脳の広い領域から入力を受け、脳と脊髄の広い領域に出力を送る。ラットでは、扁桃体に向かう群は恐怖の学習を強め、前頭前皮質に向かう別の群はその消去を強める。 & \cite{Schwarz2015LC, Uematsu2017LC, Poe2020} & 測定済み \\
3 & 持続モード：毎秒1回未満から数回の規則的なスパイクで、発火率は覚醒度とともに変わる & ラット、睡眠覚醒サイクルでの青斑核ニューロンの記録。発火率は覚醒時に最大、徐波睡眠で低く、レム睡眠でほぼ0。発火率の変化は状態の切り替わりに先行する。 & \cite{AstonJonesBloom1981} & 測定済み \\
4 & 相動モード：課題にとって重要な出来事に、ほぼ決定の瞬間に短いバースト & まれな標的を見つける課題のサル。青斑核のすべてのニューロンが標的にだけバーストで応答し、その$\sim90\ms$後、手の応答の$\sim200\ms$前に出る。成績が悪いときはバーストが弱い。二者択一の課題では、バーストは刺激より応答に正確に結びつき、正答でも誤答でも出る。 & \cite{AstonJones1994vigilance, Clayton2004LC} & 測定済み \\
5 & 瞳孔径は\nt{NORA}の不正確なテストポイント：\nt{NORA}にも\nt{ACET}にも他の領域にも従う & サル。青斑核のスパイクは、1細胞の個々のスパイクに至るまで瞳孔の散大を予測する。似ているがより遅い関係が丘と帯状皮質にもある。マウス。瞳孔のゆらぎは、皮質のノルアドレナリン作動性出力とコリン作動性出力の両方の活動に従う。 & \cite{Joshi2016, Reimer2016} & 測定済み \\
6 & ノルアドレナリンは背景活動を誘発活動より強く抑える。乗算的ゲイン~\eqref{eq:gain}はこの効果を再現するモデル & 覚醒サルの聴覚皮質、ノルアドレナリンのイオントフォレシス。ニューロンの背景活動は同種の声への応答より強く抑えられ、SN比が上がる。乗算的シグモイド~\eqref{eq:gain}はネットワークモデルからとったもので、傾きの文字どおりの乗算は測定されていない。 & \cite{Foote1975NE, ServanSchreiber1990} & 測定済み。$g$はモデル \\
7 & ゲインが$d'$を上げるのは増幅器の後のノイズに対してだけ~\eqref{eq:gainsnr}（命題~\ref{prop:gainnoise}） & 本章で導出。ネットワークモデルではゲインは信号検出を改善する。増幅器の前と後のノイズを分けてこの違いを検証した実験はない。 & 本章で導出、\cite{ServanSchreiber1990} & 理論 \\
8 & 境界$g\beta=1$は選択ネットワークを「思案中」から「決定済み」に切り替える~\eqref{eq:wtagain}（命題~\ref{prop:gainwta}、図~\ref{fig:pitchfork}） & 互いに抑制し合う二つの群の線形解析。本章で導出。脳でのこの閾値の直接測定はない。 & 本章で導出 & 理論 \\
9 & ゲインは選択の逆温度~\eqref{eq:softmax} & 増幅器の後のロジスティックノイズでは、選択確率は$T=\sigma/g$のソフトマックスになる。本章で導出。 & 本章で導出 & 理論 \\
10 & \nt{NORA}はどれだけ探索するかを決める：高い持続性活動は小さな実効$g$（探索）、決定の瞬間の相動性バーストは大きな実効$g$（決定） & ヒト。当て推量の選択の前には、既知の最良の選択の前よりベースラインの瞳孔が広い。瞳孔が広い人ほど探索の傾向が強い。ラット。ランダムな選択のモードへの移行は、前帯状皮質への\nt{NORA}入力が制御する。バーストが実効$g$を上げることは直接測定されていない。 & \cite{JepmaNieuwenhuis2011, Tervo2014stochastic, AstonJones2005} & 仮説 \\
11 & 報酬は$g$に対して逆U字に依存する。速く変わる世界では最良の$g$は小さい（命題~\ref{prop:explore}、図~\ref{fig:invertedU}） & \texttt{models/nora\_explore.py}、$4000$試行の実行$60$回。$1000$試行に1回の変化：最良$g=16$、割合$0.976$。$20$試行に1回：$g=8$、割合$0.886$。$g=0.5$では$0.77$。再計算は本章と一致した。 & \texttt{models/\allowbreak nora\_\allowbreak explore.py} & モデル \\
12 & 行動における逆U字：明瞭なバーストを伴う中程度の持続性活動で成績が最もよい & 4行と同じ課題のサル。成績のよい時期には持続性発火率が中程度で、標的へのバーストが明瞭である。持続性発火率が高いとバーストがなく、誤反応が多い。持続性活動と誘発活動の変化は成績のゆらぎに従う。 & \cite{Usher1999LC, AstonJones2005} & 測定済み \\
13 & \nt{NORA}は予期しない不確実性を、\nt{ACET}は予期された不確実性を伝える & 二種類の不確実性をもつベイズ推論の理論。引用した研究には、これらの信号を伝達物質で分ける直接の実験はない。 & \cite{YuDayan2005} & 仮説 \\
14 & 急な変化の後、人はより速く学び、瞳孔が広がる & ヒト、急な変化のある予測課題。瞳孔の短い変化は変化確率の推定に従い、ベースラインの瞳孔とともに、新しいデータが推定をどれだけ変えるか（学習率）を予測する。光や課題と無関係な瞳孔の変化もこの率を変える。ここで瞳孔がまさに\nt{NORA}を示すというのは、5行の間接的データからの推論である。 & \cite{Nassar2012} & 測定済み \\
15 & 相動性バーストは割り込みとして働く：ネットワークは状態をリセットする & \nt{NORA}のバーストがネットワークの状態をリセットする直接の実験はない。測定されているのは重要な出来事へのバーストだけである（4行）。 & \cite{BouretSara2005} & 仮説 \\
16 & 驚きによる$g$や学習率の調整は、最良の一定設定よりほとんどよくない & \texttt{models/nora\_explore.py}、各時期$1000$試行の世界（$1000$試行に1回と$20$試行に1回の変化）。最良の一定$g=8$：$0.925$。$g$の調整：最大$0.927$。学習率の調整：最大$0.926$。一定の学習率$0.4$：$0.928$。再計算は本章と一致した。 & \texttt{models/\allowbreak nora\_\allowbreak explore.py} & モデル \\
\bottomrule
\end{longtable}}
